\newpage
\begin{center}
  \textbf{Resumen}
\end{center}

\noindent
La presente monografía propone el diseño e implementación de un sistema inteligente de recuperación semántica para convocatorias públicas del Sistema de Contrataciones Estatales (SICOES) de Bolivia. El problema abordado surge de las limitaciones de la búsqueda tradicional basada en coincidencias textuales, la cual puede omitir oportunidades relevantes cuando la necesidad del usuario se expresa con términos distintos a los utilizados en el objeto de contratación.

El estudio se estructura bajo la metodología CRISP-DM y contempla las fases de extracción de datos públicos, limpieza y transformación del dataset, construcción de un corpus orientado a recuperación aumentada por generación, generación de embeddings y almacenamiento vectorial en PostgreSQL con la extensión pgvector. Sobre esta base se plantea una comparación entre búsqueda tradicional mediante SQL e \texttt{ILIKE}, búsqueda semántica por similitud vectorial y generación de respuestas asistida por un modelo de lenguaje a través de LangChain.

La propuesta tiene un enfoque aplicado y reproducible, integrando notebooks para preparación de datos, una base de datos desplegable con Docker y una interfaz en Streamlit para consulta. Los resultados obtenidos muestran que, para el corpus vigente y el conjunto curado de consultas de evaluación, la búsqueda keyword por términos constituye una línea base particularmente fuerte, mientras que la búsqueda semántica pura no supera dicho baseline. En este contexto, la estrategia híbrida y la arquitectura RAG configurable aportan una alternativa más robusta para combinar coincidencia léxica, recuperación contextual y generación asistida por modelos de lenguaje.

\bigskip
\noindent \textbf{Palabras clave:} recuperación semántica, SICOES, contrataciones públicas, pgvector, PostgreSQL, LangChain, RAG.

\newpage
\section*{Abstract}

\noindent
This monograph proposes the design and implementation of an intelligent semantic retrieval system for public procurement calls published in Bolivia's State Procurement System (SICOES). The problem addressed arises from the limitations of traditional keyword-based search, which may miss relevant opportunities when the user's need is expressed with terms different from those used in the procurement description.

The study follows the CRISP-DM methodology and includes public data extraction, dataset cleaning and transformation, construction of a retrieval-oriented corpus, embedding generation, and vector storage in PostgreSQL with the pgvector extension. Based on this architecture, the project compares traditional SQL search using \texttt{ILIKE}, semantic search through vector similarity, and answer generation assisted by a language model through LangChain.

The proposal has an applied and reproducible orientation, integrating notebooks for data preparation, a Docker-based database setup, and a Streamlit interface for consultation. The obtained results show that, for the current corpus and curated evaluation queries, token-based keyword retrieval remains a strong baseline, while pure semantic retrieval does not outperform it. In this setting, hybrid retrieval and a configurable RAG architecture provide a more robust way to combine lexical matching, contextual retrieval, and language-model-assisted answer generation.

\bigskip
\noindent \textbf{Keywords:} semantic retrieval, SICOES, public procurement, pgvector, PostgreSQL, LangChain, RAG.
